\begin{explanationbox}
\textbf{\textcolor{CExplanation}{候选横坐标不等于横轴交点}}
对 $f(x)=x^3-x-1$，首项系数为 $1$，常数项为 $-1$，所以有理根候选只有 $\pm1$。图中空心圆标记这两个候选横坐标，实心点标记实际函数值。
\begin{center}
\begin{tikzpicture}[font=\small]
\begin{axis}[width=12cm,height=7cm,axis lines=middle,
 xmin=-1.5,xmax=2.2,ymin=-3.4,ymax=6.4,
 xtick={-1,1,2},xticklabel style={anchor=south,yshift=6pt},
 ytick={-1,5},xlabel={$x$},ylabel={$y$},
 clip=false]
\addplot[CStatement,very thick,domain=-1.45:2.05,samples=151] {x^3-x-1};
\addplot[only marks,mark=o,mark size=3pt,CTrap] coordinates {(-1,0) (1,0)};
\addplot[only marks,CStatement] coordinates {(-1,-1) (1,-1) (2,5)};
\node[anchor=north] at (axis cs:-0.9,-2.4) {$(-1,-1)$};
\node[anchor=north] at (axis cs:1,-1.15) {$(1,-1)$};
\node[anchor=east] at (axis cs:1.95,5.2) {$(2,5)$};
\addplot[only marks,CExample,mark size=3pt] coordinates {(1.324717957,0)};
\node[anchor=east,text=CExample] at (axis cs:1.1,2.4) {无理数零点};
\draw[CExample] (axis cs:1.324717957,0.15)--(axis cs:1.05,1.95);
\node at (axis cs:-0.65,4.8) {$y=x^3-x-1$};
\end{axis}
\end{tikzpicture}
\end{center}

\textbf{看图：} 两个候选位置都没有落在曲线上：$f(-1)=f(1)=-1\ne0$，故没有有理根。可是 $f(1)<0<f(2)$，由连续性，$(1,2)$ 内仍有实根，且只能是无理数。绿色点由原函数数值求零点定位，横坐标约为 $1.324718$（四舍五入至六位小数）；无理性的判定依靠前面的候选排除。

\medskip
\textbf{\textcolor{CExplanation}{把图像观察变成操作}}
\begin{enumerate}[leftmargin=*,itemsep=0.5em]
\item \textbf{先限：} 把系数化为整数；常数项为零先提 $x$。其余情形用 $p\mid a_0$、$q\mid a_n$ 列不同的最简分数。
\item \textbf{再验：} 先试 $\pm1$ 等简单数，再试其他候选；必须计算 $f(p/q)$。
\item \textbf{后除：} 验出根才提取 $qx-p$。继续分析商式，最后将各因式乘回核验，防止遗漏重根。
\end{enumerate}
最简分数不可省：$2/4$ 应写为 $1/2$，再取 $p=1,q=2$。\textbf{没有有理根}只排除有理横坐标，不排除无理实根，也不排除二次及更高次因式。
\end{explanationbox}
\clearpage
